\scanpage{16}
\mspage{7}
\uline{Proposition} Si $(G=(S,\widetilde A,\sigma),I,\uline{g}:S\to\mathbb N)$ est un MD-graphe, son type $(g,\nu)$ est anabélien, i.e. $2g+\nu\geq3$.

Si on avait $g=1$, $\nu=0$, alors la relation
\[
g=1=\sum g_\alpha+h_1
\]
montre que ou bien tous les $g_\alpha$ sont nuls, et $h_1=1\deleted{\unsure{+1}}$, ou bien tous les $g_\alpha$ sauf un \inserted{$g_{\alpha_0}$} sont nuls, et $g_{\alpha_0}=1$, $h_1=0$ \inserted{(i.e. $G$ est un arbre)}. \inserted{Dans le premier cas,} La relation $s_0-s_1=h_0-h_1=1-1=0$ montre qu'il y a autant de sommets que d'arêtes, dans le deuxième, $s_0-s_1=1-0=1$ il y a un sommet de plus que d'arêtes.

Mais la relation $\nu=0$ implique que les sommets de genre $0$ sont d'ordre $\geq3$. \inserted{$G$ n'est pas réduit à un sommet, puisqu'il est stable de type $(1,0)$.} Mais un \deleted{graphe top. compact \unsure{ou}} \inserted{arbre qui \unsure{n'est réduit} à un seul sommet} a au moins deux bouts, dont l'un au moins sera de genre $0$, ce qui élimine le cas 2°. Donc tous les sommets sont de genre nuls \inserted{(donc aucun sommet n'est un bout)} et $h_1=\unsure{1}$. Mais \deleted{\unsure{comme}} un graphe topologique $0$-connexe avec $h_1=1$ qui n'a pas de \inserted{pt.} bord topologique est homéomorphe au cercle, donc s'il est réalisé par un graphe combinatoire, tous les sommets de celui-ci sont d'ordre $2$, ce qui contredit qu'ils soient d'ordre trois.
\mspage{15}
